\section{Compression of bounded quadratic-tilt towers}\label{sec:critical-compression}
This section uses $\phi=\tanh$. Write $\|A\|_{\max}=\max_{ij}|A_{ij}|$ and
$\|A\|_{1,\mathrm{ent}}=\sum_{ij}|A_{ij}|$ for matrices on the finite input
set. These are different norms; their dual pairing is what avoids a factor
of $2^d$ below.
For a positive semidefinite covariance $Q$ and a symmetric matrix $\Gamma$,
let
\[
 \mathsf F_\Gamma(Q)=
 \frac{\mathbb E_{Z\sim\mathcal N(0,Q)}
 [\phi(Z)\phi(Z)^\top\exp\{\phi(Z)^\top\Gamma\phi(Z)\}]}
 {\mathbb E_{Z\sim\mathcal N(0,Q)}
 [\exp\{\phi(Z)^\top\Gamma\phi(Z)\}]}.
\]
All activations are coordinatewise. The definition is valid for singular $Q$.

\begin{lemma}[A dimension-independent covariance comparison]\label{lem:tilted-covariance}
If $g=\|\Gamma\|_{1,\mathrm{ent}}$, then for all positive semidefinite
$Q,Q'$ of the same size,
\[
 \|\mathsf F_\Gamma(Q)-\mathsf F_\Gamma(Q')\|_{\max}
 \leq 8(1+g)^2\|Q-Q'\|_{\max}.
\]
\end{lemma}
\begin{proof}
Put $V(z)=\phi(z)^\top\Gamma\phi(z)$ and
$Q_t=(1-t)Q+tQ'$. For a twice differentiable bounded test function $F$, let
$\mathbb E_t$ denote expectation under the Gaussian of covariance $Q_t$
reweighted by $e^V$. Gaussian differentiation and the quotient rule give
\begin{align*}
 \frac{d}{dt}\mathbb E_tF
 =\frac12\sum_{a,b}(Q'-Q)_{ab}\bigl\{
 &\mathbb E_t[\partial_{ab}F+
       \partial_aF\partial_bV+\partial_bF\partial_aV]\\
 &+\operatorname{Cov}_t(F,\partial_{ab}V+
                                   \partial_aV\partial_bV)\bigr\}.
\end{align*}
This formula has no reciprocal partition-function bound: all terms are
expectations under a probability measure. Define
$M_1(H)=\sup_z\sum_a|\partial_aH(z)|$ and
$M_2(H)=\sup_z\sum_{a,b}|\partial_{ab}H(z)|$.
Because $|\phi|,|\phi'|\leq1$ and $|\phi''|\leq2$,
\[
 M_1(V)\leq2g,\qquad M_2(V)\leq6g.
\]
For $F(z)=\phi(z_i)\phi(z_j)$, including $i=j$,
$\|F\|_\infty\leq1$, $M_1(F)\leq2$, and $M_2(F)\leq6$.
The absolute derivative is therefore at most
\[
 \frac12\|Q-Q'\|_{\max}
 \{6+8g+2(6g+4g^2)\}
 =(3+10g+4g^2)\|Q-Q'\|_{\max}.
\]
Integrating in $t$ proves the assertion, with the displayed slightly larger
constant. One may initially replace both covariances by $Q+\epsilon I$ and
$Q'+\epsilon I$. The functions and all derivatives used above are bounded,
and the partition functions lie between $e^{-g}$ and $e^g$, so weak Gaussian
convergence lets $\epsilon\downarrow0$. This also proves the singular case.
\end{proof}

\begin{proposition}[Finite realization of a bounded-tilt tower]\label{prop:tilt-compress}
Fix $L\geq2$, $K\geq1$, and $G\geq1$. Let $Q_0$ be the covariance of the
Gaussian first-layer field $w\cdot\bar x$, with
$w\sim\mathcal N(0,I_{d+1}/d)$, indexed by $x\in\mathcal X$.
Suppose symmetric matrices $\Gamma_1,\ldots,\Gamma_{L-1}$ and a vector
$\beta\in\mathbb R^{\mathcal X}$ satisfy
\[
 \|\Gamma_s\|_{1,\mathrm{ent}}\leq G,\qquad
 \|\beta\|_1\leq4K,
 \qquad Q_s=\mathsf F_{\Gamma_s}(Q_{s-1}).
\]
Then the function $z(x)=(Q_{L-1}\beta)_x$ can be uniformly approximated to
accuracy $1/8$ by a depth-$L$ network whose width and weights are bounded by
$\operatorname{poly}_L(d,K,G)$.
If $f(x)z(x)\geq1/2$ everywhere, a factor-two increase in the finite
network's readout implements $f$ with margin at least $3/4$.
\end{proposition}
\begin{proof}
Let $b=8(1+G)^2$ and
\[
 \eta=[32K(L-1)b^{L-2}]^{-1}.
\]
Choose the common width $n$ to be the least integer satisfying
\[
 n\geq2\eta^{-2}\log\bigl(8(L-1)\,2^{2d}\bigr).
\]
The construction is random only to prove existence.
The matrices $\Gamma_s$ may be indexed by the entire cube; they select the
sampling laws, but are not stored or evaluated by the resulting finite
network. No efficient sampling or construction algorithm is claimed.
At the first layer, independently sample $n$ Gaussian weight vectors
$w_i\sim\mathcal N(0,I_{d+1}/d)$ tilted by
\[
 \exp\{h(w_i)^\top\Gamma_1h(w_i)\},\qquad
 h(w_i)_x=\phi(w_i\cdot\bar x).
\]
Their activation vectors have exactly the law defining
$Q_1=\mathsf F_{\Gamma_1}(Q_0)$.
At each subsequent hidden layer, conditional on the already sampled previous
activation vectors $h_1,\ldots,h_n$, independently sample $n$ standard
Gaussian row vectors $\xi_j\in\mathbb R^n$ tilted by
\[
 \exp\{\phi(Z_j)^\top\Gamma_s\phi(Z_j)\},
 \qquad Z_j(x)=\frac1{\sqrt n}\sum_{i=1}^n\xi_{ji}h_i(x).
\]
Use the actual network weights $\xi_{ji}/\sqrt n$.
If $\widehat Q_{s-1}=n^{-1}\sum_i h_ih_i^\top$ is the preceding empirical
activation covariance, then, conditionally, $Z_j$ before tilting is Gaussian
with exactly that covariance. Thus the new activation vectors are
conditionally independent, with second moment
$\mathsf F_{\Gamma_s}(\widehat Q_{s-1})$.

Every covariance entry is the empirical mean of $n$ variables in $[-1,1]$.
Hoeffding's inequality, a union bound over $2^{2d}$ entries, and conditioning
at successive layers show, with probability at least $3/4$, that
\[
 \|\widehat Q_s-
          \mathsf F_{\Gamma_s}(\widehat Q_{s-1})\|_{\max}\leq\eta
 \quad(1\leq s<L),
\]
where $\widehat Q_0=Q_0$.
Lemma~\ref{lem:tilted-covariance} then implies
\[
 \|\widehat Q_{L-1}-Q_{L-1}\|_{\max}
 \leq (L-1)b^{L-2}\eta.
\]
For each last-layer activation vector $h_i$, choose the deterministic
readout weight $v_i=(\beta\cdot h_i)/n$. The resulting output is
$\widehat z(x)=(\widehat Q_{L-1}\beta)_x$, and therefore
\[
 \|\widehat z-z\|_\infty
 \leq4K(L-1)b^{L-2}\eta=1/8.
\]
Although $\beta$ is indexed by the entire input cube, its only role in the
finite network is to select these $n$ real readout weights; each satisfies
$|v_i|\leq4K/n$. It is not an additional computational input to the network.

It remains to bound the sampled weights without losing polynomial size.
For every sampled row, its tilted density relative to its Gaussian reference
is at most $e^{2G}$, since the quadratic tilt lies in $[-G,G]$.
Let $D_n=n(d+1)+(L-2)n^2$ be the total number of incoming scalar coordinates,
and choose
\[
 B=\bigl[4G+2\log(8\max\{D_n,1\})\bigr]^{1/2}.
\]
Each raw incoming coordinate has reference variance at most one, so, also
conditionally on all preceding layers, its probability of exceeding $B$ in
absolute value is at most $2e^{2G-B^2/2}$.
A union bound gives probability at most $1/4$ that any incoming coordinate
exceeds $B$. The actual internal weights divide the raw coordinates by
$\sqrt n$ and so satisfy the same bound. Intersecting this event with the
covariance-approximation event leaves probability at least $1/2$.
Fix a realization in the intersection. Both $n$ and
$\max\{B,8K/n\}$ after the factor-two readout increase are bounded by
polynomials in $d,K,G$ at fixed depth. The stated margin follows from
$2(1/2-1/8)=3/4$.
\end{proof}

\begin{lemma}[Counting finite realizations]\label{lem:tilt-count}
Fix $L$, width $n$, and an entrywise bound $B\geq1$. The Boolean functions
having depth-$L$, width-$n$ implementations with margin at least $3/4$ and
weights bounded by $B$ number at most
\[
 \exp\{P\log(3+16B\Lambda)\},\qquad
 P=n(d+1)+(L-2)n^2+n,
 \quad \Lambda=(d+Ln)[1+n(B+1)]^{L-1}.
\]
\end{lemma}
\begin{proof}
Round every weight to a multiple of $(8\Lambda)^{-1}$.
The first hidden layer's coordinatewise error is at most $(d+1)/(8\Lambda)$.
Every subsequent hidden layer has error at most $n(B+1)$ times the preceding
error plus $n/(8\Lambda)$; the same estimate applies at the linear readout.
Induction bounds the final error by $1/8$. Each rounded coordinate has at
most $3+16B\Lambda$ possible values, so the $P$-parameter grid has the stated
size. Rounding preserves every represented Boolean sign function.
\end{proof}
